\chapter{光与物质的相互作用：光的量子性}
\problem{7-1}{有一介质，吸收系数 $\alpha=0.32\,\mathrm{cm}^{-1}$，透射光强分别为入射光强的10\%、20\%、50\%和80\%时，介质的厚度各若干？}
\answerarea
\problem{7-2}{一玻璃管长 $3.50\,\mathrm m$，内贮标准大气压下的某种气体，若这气体在此条件下的吸收系数为 $0.1650\,\mathrm m^{-1}$，求透射光强的百分比。}
\answerarea
\clearpage
\problem{7-3}{一块光学玻璃对水银灯蓝、绿谱线 $\lambda=435.8\,\mathrm{nm}$ 和 $546.1\,\mathrm{nm}$ 的折射率分别为 $1.65250$ 和 $1.62450$。用此数据定出科西公式（7.11）中的 $A$、$B$ 两常量，并用它计算对钠线 $\lambda=589.3\,\mathrm{nm}$ 的折射率 $n$ 及色散率 $\mathrm dn/\mathrm d\lambda$。}
\answerarea
\problem{7-4}{利用第一章表1-3中冕牌玻璃 K9 对 F、D、C 三条谱线的折射率数据定出科西公式（7.10）中的 $A$、$B$、$C$ 三常量，用它计算该表中给出的其它波长下折射率数据，并与表中实测数值比较。\missing{原书表1-3。}}
\answerarea
\clearpage
\problem{7-5}{一棱镜顶角 $50^\circ$，设它的玻璃材料可用二常量科西公式（7.11）来描写，其中 $A=1.53974$，$B=4.6528\times10^3\,\mathrm{nm^2}$。求此棱镜对波长 $550.0\,\mathrm{nm}$ 周围最小偏向角的色散本领。}
\answerarea
\problem{7-6}{根据（7.25）式证明吸收峰的高度反比于 $\gamma_j$，半值宽度（即峰值之半处的宽度，见图7-9b）$\Delta\lambda$ 正比于 $\gamma_j$。
\begin{equation*}
n\kappa(\lambda)=\frac{Ne^2}{2\varepsilon_0m}\frac{1}{2\pi c}
\sum_j\frac{f_j\gamma_j\lambda_j^4\lambda^3}
{(2\pi c)^2(\lambda^2-\lambda_j^2)^2+\gamma_j^2\lambda_j^4\lambda^2}.
\tag{7.25}
\end{equation*}
其中 $\lambda_j$ 为第 $j$ 类振子的共振波长，$\gamma_j$ 为阻尼参数，$f_j$ 为相应的振子强度。
\bookfig[0.88]{book_7_9}{图7-9\quad 振子的色散与共振吸收曲线}}
\answerarea
\clearpage
\problem{7-7}{一块玻璃对波长 $0.070\,\mathrm{nm}$ 的 X 射线的折射率比1小 $1.600\times10^{-6}$。求 X 射线能在此玻璃外表面发生全反射的最大掠射角。}
\answerarea
\problem{7-8}{估计一下铜的等离子体振荡角频率 $\omega_p$ 的数量级。}
\answerarea
\clearpage
\problem{7-9}{求习题7-4中冕牌玻璃 K9 对 D 双线的群速。\missing{原书表1-3。}}
\answerarea
\problem{7-10}{试计算下列各情况下的群速：
\begin{enumerate}\item $v_p=v_0$（常量）（无色散介质，如空气中的声波）。\item $v_p=\sqrt{\dfrac{\lambda}{2\pi}\left(g+\dfrac{4\pi^2T}{\lambda^2\rho}\right)}$（水面波，$g$ 为重力加速度，$T$ 为表面张力，$\rho$ 为液体的密度）。\item $n$ 满足正常色散的科西公式（7.11）。\item $\omega^2=\omega_c^2+c^2k^2$（波导中的电磁波，$\omega_c$ 为截止角频率）。\end{enumerate}}
\answerarea
\clearpage
\problem{7-11}{摄影者知道用橙黄色滤色镜拍摄天空时，在黑白照片中可增加蓝天和白云的对比。设照相机的镜头和底片的灵敏度将光谱范围限制在 $390.0\,\mathrm{nm}$ 到 $620.0\,\mathrm{nm}$ 之间，并设天空散射光在此范围内可看成是常量。若滤色镜把波长在 $550.0\,\mathrm{nm}$ 以下的光全部吸收，天空的散射光被它去掉了百分之几？}
\answerarea
\problem{7-12}{苯（$\mathrm{C_6H_6}$）的拉曼散射中较强的谱线与入射光的波数差为607、992、1178、1586、3047、$3062\,\mathrm{cm}^{-1}$。今以氩离子激光（$\lambda=488.0\,\mathrm{nm}$）入射，计算各斯托克斯和反斯托克斯光谱的波长。}
\answerarea
\clearpage
\problem{7-13}{设一个两能级系统能级差 $E_2-E_1=0.01\,\mathrm{eV}$。
\begin{enumerate}\item 分别求 $T=10^2\,\mathrm K$、$10^4\,\mathrm K$、$10^5\,\mathrm K$、$10^8\,\mathrm K$ 时粒子数 $N_2$ 与 $N_1$ 之比；\item $N_2=N_1$ 的状态相当于多高的温度？\item 粒子数发生反转的状态相当于怎样的温度？\item 我们姑且引入“负温度”的概念来描述粒子数反转的状态，你觉得 $T=-10^8\,\mathrm K$ 和 $T=+10^8\,\mathrm K$ 两个温度中哪一个更高？\end{enumerate}}
\answerarea
\clearpage
